\subsection{Review: Adding and Subtracting Fractions} 
We can add fractions with the same \makeblank[1.5in].
\begin{procedure}{Adding or Subtracting Fractions with Like Denominator}
If fractions have the same denominator, we can add or subtract them by adding or subtracting the numerators.
\end{procedure}
\begin{example}
Add or subtract. Simplify your answer as much as possible.
\begin{lpic}{}{4}
\foreach \y/\l/\na/\nb/\d/\s in {1/a/5/7/9/+,3/b/9/3/8/-} {
	\fractextright{1}{-\y}{\l. };
	\fractextleft{1}{-\y}{$\dfrac{\na}{\d}\s\dfrac{\nb}{\d}={}$};
}
\end{lpic}
\end{example}
Otherwise, we have to \makeblank[1.25in] them to have the same denominator.
\begin{procedure}{Adding or Subtracting Fractions}
\begin{enumerate*}
\item Find the least common multiple of the denominators (LCD).
\item Write equivalent fractions with that denominator.
\item Add the new fractions (that have the same denominator).
\item Simplify, if necessary.
\end{enumerate*}
\end{procedure}
There are many methods we can use to find a common denominator with numbers. In this example, we'll use the method we'll use for rational expressions.
\begin{example}
Add $\dfrac{5}{6}+\dfrac{7}{10}$.
\begin{lpic}{}{12}
\foreach \y/\n/\text in {1/1/Find the LCD by looking at factors.,
												4/2/Rewrite the fractions with the LCD.,
												7/3/Add.,
												10/4/Simplify if necessary.} {
	\regtextright{1}{-\y}{\n. };
	\regtextleft{1}{-\y}{\text};
}
\foreach \x/\n/\d in {1/5/6,\value{cols}/7/10} {
	\regtextleft{\x}{-2}{$\d={}$};
	\fractextleft{\x}{-5}{$\dfrac{\n}{\d}={}$};
}
\regtextleft{1}{-3}{$\mathrm{LCD}={}$};
\fractextleft{1}{-8}{$\dfrac{5}{6}+\dfrac{7}{10}={}$};
\end{lpic}
\end{example}